\documentclass[11pt]{article}
\usepackage[a4paper,margin=27mm]{geometry}
\usepackage[T1]{fontenc}
\usepackage{lmodern,microtype,amsmath,amssymb,amsthm,mathtools,enumitem}
\usepackage[hidelinks,hypertexnames=false]{hyperref}
\newtheorem{theorem}{Theorem}
\newtheorem{lemma}[theorem]{Lemma}
\newtheorem{proposition}[theorem]{Proposition}
\newtheorem{corollary}[theorem]{Corollary}
\theoremstyle{definition}\newtheorem{definition}[theorem]{Definition}
\theoremstyle{remark}\newtheorem{remark}[theorem]{Remark}
\newcommand{\R}{\mathbb R}
\newcommand{\Dscr}{\mathcal D}
\newcommand{\supp}{\operatorname{supp}}
\newcommand{\debt}{\mathrm{debt}}
\newcommand{\cred}{\mathrm{credit}}
\newcommand{\sil}{\mathrm{silence}}
\newcommand{\Band}{\mathrm{B}}
\title{Closed Choice Forces Fineness}
\author{}\date{10 October 2026}
\begin{document}\maketitle

\noindent
This note continues IC (\emph{Independence Chains and Two Additive Worlds})
and DT (\emph{Dense Thresholds and Debt Reduction}). It settles the case of
two additive worlds studied there, provided the third world's ranking is
complete. It also reduces the conjecture for complete rankings to a property
of edge selections.

\begin{enumerate}[leftmargin=*,itemsep=3pt]
\item \emph{Main theorem} (Theorem~\ref{thm:two}). Let $n=3$ and assume (G)
and (S). Suppose two worlds have one-sided weak extensions of CBOI's finitely
additive form; ordinary additive losses are a special case. Suppose the third
world's ranking is complete and satisfies (IB). Then $V$ is dominated.
\item In particular, the sub-case of IC and DT is settled whenever world~$1$
is complete. This includes DT's residual configuration
(Corollary~\ref{cor:SH}). Of the chain arguments of IC and DT, only IC's
chain lemma is used (through Proposition~\ref{prop:FA}).
\item \emph{Reduction} (Theorem~\ref{thm:reduction}). Let $n\ge3$, let every
ranking be complete with (G) and (IB), and assume (S). Suppose every
selection on every edge contains an endpoint of that edge. Then every
coherent model whose credal set contains three vertices is dominated.
\item \emph{Consequence} (Corollary~\ref{cor:shape}). Let $n=3$, with
complete rankings, and suppose $V$ is undominated. Then at most one world has
such a weak extension, so at least two worlds carry leverage sequences in
EDL's sense. Moreover, some edge has a selection that omits both endpoints.
\end{enumerate}
The new ingredient is the following. By (IB), the comparisons at a single
world of pure debt sets, and separately of pure credit sets, are comparative
probabilities in de Finetti's sense. The closed sections of (S) at the edge
selections make them \emph{fine}. The Savage--Wakker theorem then gives
finitely additive probabilities that almost agree with them. One exchange
rate between debt and credit assembles these into a weak extension of the
form required by CBOI's theorem.

\emph{What remains open.} First, the case where the third world's preorder
is incomplete (Remark~\ref{rem:complete}). Second, the endpoint property
without any additive world, which is now the whole conjecture for complete
rankings. Section~\ref{sec:remarks} also records a referee report on DT. It
found one gap there, in the lemma on debt contraction, and a repair.

Two independent referees checked the argument: the first the core argument
for two exactly additive worlds, the second this note in full. Their
corrections are incorporated. The proofs of Lemma~\ref{lem:rates} and
Theorem~\ref{thm:extension} are the first referee's replacements for proofs
that had gaps.

\section{Setting}\label{sec:setting}

We use the setting of v3, CC, CBOI, EDL, IC and DT. Binary problems
$\{g,-g\}$ are represented by $g\in H\coloneqq\{g\in\R^n:g_1>0\}$, with
normalized Gaussian measure $\mu$. The mixed problems are
$H'\coloneqq\{g\in H:g_j<0\text{ for some }j\}$; the remaining problems in
$H$ are unanimous (up to null sets). A regular model is correct on unanimous
problems and never contradictory. On $H'$ it is an advice function with three
values. At world~$1$ these are \emph{credit} (accept $g$, rank $0$),
\emph{silence} (rank $1$) and \emph{debt} (accept $-g$, rank $2$). For a
regular $M$ write $\cred(M)$ and $\debt(M)$ for the corresponding subsets of
$H'$. For measurable $Q\subseteq H'$, $M+\debt(Q)$ is the model that agrees
with $M$ except that it gives debt on $Q$. Similarly for credit and silence.

(G) is strict guidance monotonicity and (IB) independence under replacement
of shared advice. (S) supplies coherent undominated models $C_p$,
$p\in\Delta$, competitive on $\supp p$, with closed sections
$A_i(X)=\{p:C_p\preceq_iX\}$. (S) implies (LA). We use the following
consequences of (S): every model has a selected coherent weak dominator at
all worlds (CBOI, Lemma~1); and the face lemma (CC, Lemma~7), by which the
credal set of $C_q$ lies in $\Delta_{\supp q}$. If $V$ is undominated, it is
competitive on all worlds (IC, Section~1).

For $j\ge2$ put
\[
 D_j\coloneqq\{g\in H:g_j<0\},\qquad t_j\coloneqq g_1/|g_j|\ \text{ on }D_j.
\]
Then $H'=\bigcup_{j\ge2}D_j$. The law of $t_j$ under $\mu|_{D_j}$ has a
positive density on $(0,\infty)$, and for $j\ne k$ the joint law of
$(t_j,t_k)$ on $D_j\cap D_k$ has a positive density on $(0,\infty)^2$. For
$0\le s<s'\le\infty$ put
\[
 \Band_j[s,s')\coloneqq D_j\cap\{s\le t_j<s'\},\qquad
 \Band_j(s,s']\coloneqq D_j\cap\{s<t_j\le s'\}.
\]

\paragraph{Losses of finitely additive form.} Following CBOI, let
$\nu_{k1},\nu_{k2},\nu_{k3}$ be finite nonnegative finitely additive
measures on the measurable subsets of $H$ with
\begin{equation}
 \nu_{kl}(A)=0\text{ if }\mu(A)=0,\qquad\nu_{kl}(A)>0\text{ if }\mu(A)>0,
 \label{eq:positive}
\end{equation}
and put $L_k(\Dscr)\coloneqq\sum_{l=1}^3\nu_{kl}\{g:r_k(g,\Dscr)\ge l\}$.
Ordinary additive losses with strictly ordered local penalties and an
integrable envelope are the special case of densities. Regular models never
use rank~$3$, and are correct on unanimous problems. So for regular models
only $\nu_{k1},\nu_{k2}$ on $H'$ matter, and $\nu_{k3}$ and the restrictions
to unanimous problems may be replaced by $\mu$ without changing $L_k$ on
regular models. Consider the one-sided bound
\begin{equation}
 B\text{ coherent undominated},\ X\text{ regular},\ B\preceq_kX
 \ \Longrightarrow\ L_k(B)\le L_k(X).
 \tag{W$_k$}\label{eq:W}
\end{equation}
It needs no completeness of $\preceq_k$.

\begin{lemma}\label{lem:cboi}
Let $n\ge3$ and assume (G) and (S). Suppose every world $k$ has a loss $L_k$
of the above form satisfying \eqref{eq:positive} and \eqref{eq:W}. Then every
coherent model whose credal set contains three vertices is dominated.
\end{lemma}
\begin{proof}
This is CBOI's Theorem~2, in the one-sided form of EDL's Proposition~1. Its
proof uses the weak-extension property only for
pairs $(C_m,X_m)$, with $C_m$ selected and $X_m$ regular, and for pairs
$(T,X_m)$, with $T$ the coherent target assumed undominated and $X_m$
regular. Both kinds of pair are covered by \eqref{eq:W}.
\end{proof}

\section{Edge selections with the endpoint property}\label{sec:edges}

Fix $j\ge2$. A selection $C_q$ with $q$ on the edge $[e_1,e_j]$ has a credal
interval on that edge (face lemma). Off $D_j$ such a model gives credit,
since there $g_1>0$ and $g_j\ge0$. On $D_j$ it is determined by two
thresholds in $t_j$. The two families with an endpoint are:
\begin{itemize}[leftmargin=*,itemsep=1pt]
\item $E^j_1(a)$, $a\in[0,\infty]$, whose credal set contains $e_1$: credit
on $\Band_j(a,\infty)$, silence on the rest of $D_j$;
\item $E_j(b)$, $b\in[0,\infty]$, whose credal set contains $e_j$: debt on
$\Band_j[0,b)$, silence on the rest of $D_j$.
\end{itemize}
Then $E^j_1(0)=D_1$, $E_j(\infty)=D_j$ and
$E^j_1(\infty)=E_j(0)=V_{1j}\coloneqq\Dscr_{[e_1,e_j]}$. Index both families
by one parameter $s\in[-\infty,\infty]$:
\[
 M^j_s\coloneqq E^j_1(-1/s)\ (s<0),\qquad M^j_{-\infty}\coloneqq D_1,\qquad
 M^j_s\coloneqq E_j(s)\ (s\ge0).
\]
If $s<s'$, the models $M^j_s$ and $M^j_{s'}$ differ only on $D_j$. There the
world-1 rank of $M^j_s$ is pointwise no larger, and strictly smaller on a
non-null band. On $D_j$ the world-$j$ ranks are reversed: accepting $g$ is
bad at $j$ and accepting $-g$ is good. By strict (G),
\begin{equation}
 s<s'\ \Longrightarrow\ M^j_s\prec_1M^j_{s'}\quad\text{and}\quad
 M^j_{s'}\prec_jM^j_s.
 \label{eq:order}
\end{equation}

\begin{definition}\label{def:EP}
The \emph{endpoint property} (EP$_{ij}$) holds if every selection $C_q$ with
$q\in[e_i,e_j]$ has $e_i$ or $e_j$ in its credal set.
\end{definition}

\begin{lemma}\label{lem:edge}
Assume (G) at worlds $1$ and $j$, (S), and (EP$_{1j}$). Then there is a
continuous map $s\colon[e_1,e_j]\to[-\infty,\infty]$ with $C_q=M^j_{s(q)}$,
$s(e_1)=-\infty$ and $s(e_j)=\infty$. Every $M^j_s$ is selected. If $s$ is
finite, it is selected only at points of the open edge.
\end{lemma}
\begin{proof}
By (EP$_{1j}$) and the face lemma, each $C_q$ on the edge is some $E^j_1(a)$
or $E_j(b)$, hence some $M^j_s$. The parameter is unique, since different
parameters give models that differ on a non-null band. By the face lemma,
$C_{e_1}$ has credal set $\{e_1\}$, so $C_{e_1}=D_1$; likewise $C_{e_j}=D_j$.

For every $s$ we claim
\[
 A_1(M^j_s)\cap[e_1,e_j]=\{q:s(q)\le s\},\qquad
 A_j(M^j_s)\cap[e_1,e_j]=\{q:s(q)\ge s\}.
\]
If $s(q)\le s$, then $C_q\preceq_1M^j_s$ by (G). If $s(q)>s$, then
$M^j_s\prec_1C_q$ by \eqref{eq:order}, so $C_q\not\preceq_1M^j_s$. The second
identity is proved the same way. Both sets are closed by (S). So $s$ is lower
and upper semicontinuous, hence continuous into the extended line. By the
intermediate value theorem every $s\in[-\infty,\infty]$ is attained. If $s$
is finite, then $M^j_s\notin\{D_1,D_j\}$, so it is not selected at $e_1$ or
$e_j$.
\end{proof}

\begin{lemma}\label{lem:approach}
Under the hypotheses of Lemma~\ref{lem:edge}, let $s^*\in[-\infty,\infty]$
and let $s_n\to s^*$, with all $s_n<s^*$ or all $s_n>s^*$. Then there is
$q^*$ with $s(q^*)=s^*$ such that every neighbourhood $N$ of $q^*$ has the
following property: for some $n_0$, every $s$ between $s_{n_0}$ and $s^*$
(inclusive) equals $s(q)$ for some $q\in N$.
\end{lemma}
\begin{proof}
Choose $q_n$ with $s(q_n)=s_n$ and pass to a convergent subsequence
$q_n\to q^*$. By continuity, $s(q^*)=s^*$. A neighbourhood $N$ contains an
interval $I\ni q^*$ of the edge. Its image $s(I)$ is an interval containing
$s^*$ and $s_n$ for all large $n$ in the subsequence.
\end{proof}

\section{Band lemmas at the free world}\label{sec:bands}

From now on until Section~\ref{sec:main}, world~$1$ is the \emph{free world}.
Its ranking is \emph{complete} and satisfies (G) and (IB). For measurable
$B,C\subseteq H'$ define
\begin{gather*}
 B\preceq_dC\ :\Longleftrightarrow\ V+\debt(B)\preceq_1V+\debt(C),\\
 B\preceq_cC\ :\Longleftrightarrow\ V+\cred(C)\preceq_1V+\cred(B).
\end{gather*}
So $B\preceq_dC$ says that debt on $B$ is no worse than debt on $C$, and
$B\preceq_cC$ says that credit on $B$ is worth no more than credit on $C$.
By (IB), $V$ may be replaced by any regular model silent on $B\cup C$: the
two pairs share their advice outside $B\triangle C$ and agree on it. Strict
relations are written $\prec_d$, $\prec_c$. By (G), $B\subseteq C$ implies
$B\preceq_dC$ and $B\preceq_cC$.

Two consequences of (S) are used at an edge $[e_1,e_j]$.
\begin{enumerate}[label=(T\arabic*),leftmargin=*,itemsep=2pt]
\item If $Y$ is regular and $Y\prec_1C_{q^*}$, then $Y\prec_1C_q$ for all
$q$ near $q^*$. Indeed $q^*\notin A_1(Y)$, a closed set; so near $q^*$,
$C_q\not\preceq_1Y$, and completeness gives $Y\prec_1C_q$.
\item If $F$ is regular, $q^*\in[e_1,e_j]$ and $F\prec_jC_{q^*}$, then
$C_q\preceq_1F$ for all $q$ in the open edge near $q^*$. Indeed near $q^*$,
$C_q\not\preceq_jF$ since $A_j(F)$ is closed, and $C_q$ is competitive on
$\{1,j\}$.
\end{enumerate}
Only (T1) uses completeness.

In Lemmas~\ref{lem:debtbands} and~\ref{lem:creditbands} assume (G) at
worlds $1$ and $j$, (S) and (EP$_{1j}$), and write $t=t_j$ and
$\Band=\Band_j$.

\begin{lemma}\label{lem:debtbands}
Let $A\subseteq D_j$ be non-null.
\begin{enumerate}[label=(\alph*)]
\item If $A\subseteq\{t<a\}$ and $a<b<\infty$, there is $b'\in[a,b)$ with
$\Band[s,b)\prec_dA$ for all $s\in[b',b)$.
\item If $A\subseteq\{t\ge c\}$ and $0\le b<c$, there is $b'\in(b,c]$ with
$\Band[b,s)\preceq_dA$ for all $s\in(b,b']$.
\item If $A\subseteq\{t<a\}$, there is $s>a$ with $\Band[s,\infty)\prec_dA$.
\end{enumerate}
\end{lemma}
\begin{proof}
(a) Apply Lemma~\ref{lem:approach} with $s^*=b$ and $s_n\uparrow b$,
$s_n\ge a$. It gives $q^*$ with $C_{q^*}=E_j(b)$. Let $Y$ be $E_j(b)$ with
its debt on $A$ silenced; then $Y\prec_1E_j(b)$ by strict (G). By (T1) and
Lemma~\ref{lem:approach}, there is $n_0$ with $Y\prec_1E_j(s)$ for every
$s\in[s_{n_0},b)$. The two models differ exactly on $A$, where $Y$ is silent
and $E_j(s)$ has debt, and on $\Band[s,b)$, where $Y$ has debt and $E_j(s)$
is silent. By (IB), $\Band[s,b)\prec_dA$. Put $b'\coloneqq s_{n_0}$.

(b) Apply Lemma~\ref{lem:approach} with $s^*=b$ and $s_n\downarrow b$,
$s_n<c$. It gives $q^*$ with $C_{q^*}=E_j(b)$, in the open edge. Let
$F\coloneqq E_j(b)+\debt(A)$; $E_j(b)$ is silent on $A$. Debt on $D_j$ is
good at world $j$, so $F\prec_jE_j(b)$. By (T2) and
Lemma~\ref{lem:approach}, there is $n_0$ with $E_j(s)\preceq_1F$ for every
$s\in(b,s_{n_0}]$; these indices lie in the open edge since $s$ is finite.
The two models differ exactly on $\Band[b,s)$, where $E_j(s)$ has debt, and
on $A\subseteq\{t\ge s\}$, where $F$ has debt. So $\Band[b,s)\preceq_dA$. Put
$b'\coloneqq s_{n_0}$.

(c) Apply Lemma~\ref{lem:approach} with $s^*=\infty$ and $s_n\uparrow\infty$,
$s_n>a$; then $C_{q^*}=D_j$. Let $Y$ be $D_j$ with its debt on $A$ silenced,
so $Y\prec_1D_j$. By (T1), $Y\prec_1E_j(s)$ for some $s>a$. These differ on
$A$ and on $\Band[s,\infty)$, so $\Band[s,\infty)\prec_dA$.
\end{proof}

\begin{lemma}\label{lem:creditbands}
Let $T\subseteq D_j$ be non-null.
\begin{enumerate}[label=(\alph*)]
\item If $T\subseteq\{t<a\}$ and $a<a'<\infty$, there is $a''\in[a,a')$ with
$\Band(s,a']\prec_cT$ for all $s\in[a'',a')$.
\item If $T\subseteq\{t>c\}$ and $0\le a'<c$, there is $a''\in(a',c]$ with
$\Band(a',s]\preceq_cT$ for all $s\in(a',a'']$.
\item If $T\subseteq\{t<a\}$, there is $s>a$ with $\Band(s,\infty)\prec_cT$.
\end{enumerate}
\end{lemma}
\begin{proof}
(a) $E^j_1(s)$ with $s<a'$ has parameter $-1/s<-1/a'$. Apply
Lemma~\ref{lem:approach} with $s^*=-1/a'$ from below, using only $s\ge a$.
Let $Y\coloneqq E^j_1(a')+\cred(T)$; $E^j_1(a')$ is silent on $T$. Then
$Y\prec_1E^j_1(a')$. By (T1), $Y\prec_1E^j_1(s)$ for all $s\in[a'',a')$, for
some $a''$. The two models differ on $T$, where $Y$ has credit, and on
$\Band(s,a']$, where $E^j_1(s)$ has credit. So $V+\cred(T)\prec_1
V+\cred(\Band(s,a'])$, that is, $\Band(s,a']\prec_cT$.

(b) Apply Lemma~\ref{lem:approach} with $s^*=-1/a'$ (or $-\infty$ if $a'=0$)
from above. Then $C_{q^*}=E^j_1(a')$, which is $D_1$ if $a'=0$. Let $F$ be
$E^j_1(a')$ with its credit on $T$ silenced. Credit on $D_j$ is bad at world
$j$, so $F\prec_jE^j_1(a')$. By (T2), $E^j_1(s)\preceq_1F$ for all
$s\in(a',a'']$, for some $a''$; these indices lie in the open edge. The
models differ on $\Band(a',s]$, where $F$ has credit, and on
$T\subseteq\{t>s\}$, where $E^j_1(s)$ has credit. So $\Band(a',s]\preceq_cT$.

(c) Apply Lemma~\ref{lem:approach} with $s^*=0$ from below, so
$C_{q^*}=V_{1j}$. Let $Y\coloneqq V_{1j}+\cred(T)$, so $Y\prec_1V_{1j}$. By
(T1), $Y\prec_1E^j_1(s)$ for some $s>a$. These differ on $T$ and on
$\Band(s,\infty)$.
\end{proof}

\section{Fineness}\label{sec:fine}

\begin{theorem}\label{thm:fine}
Assume (G) at every world and (S). Let world~$1$ be complete with (IB), and
assume (EP$_{1j}$) for every $j\ge2$. Then for every non-null $A\subseteq H'$
there is a finite partition of $H'$ into measurable sets each $\preceq_dA$.
The same holds for $\preceq_c$.
\end{theorem}
\begin{proof}
We give the proof for debt; the credit proof is parallel, with
Lemma~\ref{lem:creditbands} and half-open bands $(\cdot,\cdot]$ in place of
Lemma~\ref{lem:debtbands} (details at the end). Null sets are absorbed into
any piece.

\emph{Step 0.} Since $H'=\bigcup_jD_j$, choose $j$ with $A\cap D_j$ non-null,
and $0<a_1<a_2<\infty$ such that $A_0\coloneqq A\cap D_j\cap\{a_1\le t_j\le
a_2\}$ is non-null. Then $A_0\preceq_dA$.

\emph{Step 1: one edge.} We claim: for every $j$ and every non-null
$A_0\subseteq D_j\cap\{t_j\ge a_1\}$, $a_1>0$, there are $s_0\in(0,a_1)$ and
a finite partition of $D_j$ into sets each $\preceq_dA_0$, one of which is
$\Band_j[0,s_0)$. Write $t=t_j$, $\Band=\Band_j$.

By Lemma~\ref{lem:debtbands}(b) with $b=0$, $c=a_1$, there is $s_0\in(0,a_1)$
with $\Band[0,s_0)\preceq_dA_0$. Put $A'\coloneqq\Band[0,s_0/2)$, so
$A'\preceq_dA_0$. By Lemma~\ref{lem:debtbands}(c) with $a=s_0$ there is
$s_1>s_0$ with $\Band[s_1,\infty)\prec_dA'$. Put $w\coloneqq s_1+1$. By
Lemma~\ref{lem:debtbands}(a), with $a=s_0/2$, there is $w''\in(s_1,w)$ with
$A_{\mathrm{up}}\coloneqq\Band[w'',w)\prec_dA'$, and
$A_{\mathrm{up}}\subseteq\{t\ge w''\}$.

For each $z\in[s_0,s_1]$, Lemma~\ref{lem:debtbands}(a), with yardstick $A'$
and $a=s_0/2<z$, gives $z'<z$ with $\Band[r,z)\prec_dA'$ for all
$r\in[z',z)$. Lemma~\ref{lem:debtbands}(b), with yardstick $A_{\mathrm{up}}$
and $c=w''>z$, gives $z''>z$ with $\Band[z,r)\preceq_dA_{\mathrm{up}}$ for all
$r\in(z,z'']$. Put $U_z\coloneqq(z',z'')$. Finitely many $U_z$ cover the
compact interval $[s_0,s_1]$. Let $\lambda>0$ be a Lebesgue number of this
cover, and choose $s_0=r_0<r_1<\dots<r_m=s_1$ with $r_{l+1}-r_l<\lambda$.
Each $[r_l,r_{l+1}]$ lies in some $U_z$. If $z\le r_l$, then
$\Band[r_l,r_{l+1})\subseteq\Band[z,r_{l+1})\preceq_dA_{\mathrm{up}}\prec_dA'$.
If $z\ge r_{l+1}$, then $\Band[r_l,r_{l+1})\subseteq\Band[r_l,z)\prec_dA'$. If
$r_l<z<r_{l+1}$, split $\Band[r_l,r_{l+1})$ into $\Band[r_l,z)$ and
$\Band[z,r_{l+1})$, each $\preceq_dA'$.

The sets $\Band[0,s_0)$, the pieces of $\Band[s_0,s_1)$ just obtained, and
$\Band[s_1,\infty)$ partition $D_j$ up to a null set. Each is $\preceq_dA_0$,
because $A'\preceq_dA_0$ and $\preceq_d$ is transitive.

\emph{Step 2: the other edges.} Apply Step~1 to the edge $[e_1,e_j]$ chosen
in Step~0. For $k\ne j$ the set
\[
 A_k\coloneqq\Band_j[0,s_0)\cap D_k\cap\{1\le t_k\le2\}
\]
is non-null, by the positive joint density of $(t_j,t_k)$, and
$A_k\preceq_d\Band_j[0,s_0)\preceq_dA_0$. Step~1 on the edge $[e_1,e_k]$ with
yardstick $A_k$ gives a finite partition of $D_k$ into sets
$\preceq_dA_k\preceq_dA_0$.

\emph{Step 3.} Enumerate $\{2,\dots,n\}$ as $j=j_1,j_2,\dots$. The sets
$Q\setminus\bigcup_{r'<r}D_{j_{r'}}$, with $Q$ a piece of the partition of
$D_{j_r}$, form a finite partition of $H'$. Each is a subset of a piece, so
it is $\preceq_dA_0\preceq_dA$.

\emph{Credit.} Steps 0 and 3 are unchanged. In Step~2 use
$A_k\coloneqq\Band_j(0,s_0]\cap D_k\cap\{1\le t_k\le2\}\subseteq\{t_k>1/2\}$.
In Step~1, with a yardstick $T_0\subseteq D_j\cap\{t>c\}$ (in Step~0 take
$c\coloneqq a_1/2$; in Step~2, $c\coloneqq1/2$): Lemma~\ref{lem:creditbands}(b) with
$a'=0$ gives $s_0\in(0,c)$ with $\Band(0,s_0]\preceq_cT_0$. Put
$T'\coloneqq\Band(0,s_0/2]\subseteq\{t<3s_0/4\}$. Part (c) gives $s_1>s_0$
with $\Band(s_1,\infty)\prec_cT'$. Part (a) at $a'=w\coloneqq s_1+1$ gives
$w''\in(s_1,w)$ with $T_{\mathrm{up}}\coloneqq\Band(w'',w]\prec_cT'$. For
$z\in[s_0,s_1]$, part (a) with yardstick $T'$ and $a=3s_0/4<z$ handles bands
$(r,z]$, and part (b) with yardstick $T_{\mathrm{up}}\subseteq\{t>w''\}$
handles bands $(z,r]$. The Lebesgue-number argument is the same.
\end{proof}

\section{Comparative probabilities}\label{sec:cp}

A \emph{comparative probability} on a $\sigma$-algebra $\mathcal A$ of
subsets of $\Omega$ is a complete preorder $\preceq$ on $\mathcal A$ with
$\varnothing\preceq B$ for all $B$, $\varnothing\prec\Omega$, and
$B\preceq C\iff B\cup E\preceq C\cup E$ whenever $E\cap(B\cup C)=\varnothing$.
It is \emph{fine} if for every $B\succ\varnothing$ there is a finite
partition of $\Omega$ into events each $\preceq B$. It is \emph{atomless} if
every $B\succ\varnothing$ is the union of two disjoint events each
$\succ\varnothing$. A finitely additive probability $P$ \emph{almost agrees}
with $\preceq$ if $B\preceq C$ implies $P(B)\le P(C)$.

\begin{lemma}\label{lem:cp}
Under the hypotheses of Theorem~\ref{thm:fine}, $\preceq_d$ and $\preceq_c$
are fine atomless comparative probabilities on the measurable subsets of
$H'$. Moreover $B\succ_d\varnothing$, and likewise $B\succ_c\varnothing$, if
and only if $B$ is non-null.
\end{lemma}
\begin{proof}
Completeness and transitivity come from those of $\preceq_1$, with the single
context $V$. Then $\varnothing\preceq_dB$ by (G), and strictly if $B$ is
non-null by strict (G). Null sets are equivalent to $\varnothing$ because
models are compared modulo null sets. Additivity is (IB): the pairs
$(V+\debt(B),V+\debt(C))$ and $(V+\debt(B\cup E),V+\debt(C\cup E))$ share
their advice off $B\triangle C$ and agree on it. A non-null set splits into
two non-null sets because $\mu$ is atomless. Fineness is
Theorem~\ref{thm:fine}. Credit is the same.
\end{proof}

\begin{theorem}\label{thm:SW}
A fine atomless comparative probability on an algebra of events has a unique
finitely additive probability that almost agrees with it.
\end{theorem}
This is due to Savage~\cite[\S3.3]{Savage}; Wakker~\cite{Wakker} proves it for
algebras, not only $\sigma$-algebras: fineness and atomlessness give almost
uniform partitions of every order (his Theorem~2), and these give the almost
agreeing probability (his Theorem~1 and Lemma~3). A finite partition into
events $\preceq B$ implies Wakker's form of fineness.

Let $P_d$ and $P_c$ be the probabilities given by Theorem~\ref{thm:SW} for
$\preceq_d$ and $\preceq_c$.

\begin{lemma}\label{lem:Pd}
Under the hypotheses of Theorem~\ref{thm:fine}:
\begin{enumerate}[label=(\roman*),leftmargin=*,itemsep=1pt]
\item If $P_d(B)<P_d(C)$ then $B\prec_dC$. If $P_c(B)<P_c(C)$ then
$B\prec_cC$.
\item $P_d(B)=0$ if and only if $B$ is null; likewise for $P_c$.
\item For measurable $R$, $v\in[0,P_d(R)]$ and $\varepsilon>0$ there is
$R'\subseteq R$ with $|P_d(R')-v|<\varepsilon$; likewise for $P_c$.
\item For measurable $R$, $\theta\in[0,1]$ and $\varepsilon>0$ there is
$R'\subseteq R$ with $|P_d(R')-\theta P_d(R)|\le\varepsilon$ and
$|P_c(R')-\theta P_c(R)|\le\varepsilon$.
\end{enumerate}
\end{lemma}
\begin{proof}
(i) If $C\preceq_dB$ held, almost agreement would give $P_d(C)\le P_d(B)$.
So $C\not\preceq_dB$, and completeness gives $B\prec_dC$.

(ii) A null $B$ is equivalent to $\varnothing$, so $P_d(B)\le0$. If $B$ is
non-null, a fine partition $Q_1,\dots,Q_m$ of $H'$ into events
$\preceq_dB$ gives $1=\sum_lP_d(Q_l)\le mP_d(B)$.

(iii) Given $\eta>0$, repeated halving of a non-null set gives a non-null
$C$ with $P_d(C)<\eta$. A fine partition of $H'$ into events $\preceq_dC$ has
pieces of $P_d$-mass $\le P_d(C)<\eta$. Intersect it with $R$. The partial
unions of the pieces increase from $0$ to $P_d(R)$ in steps $<\eta$.

(iv) Let $\eta>0$. Intersect a common refinement of partitions as in (iii),
for $P_d$ and for $P_c$, with $R$. This gives $R=R_1\cup\dots\cup R_m$,
disjoint, with $v_l\coloneqq(P_d(R_l),P_c(R_l))\in[0,\eta]^2$. Let
$S=(S_1,S_2)\coloneqq\sum_lv_l$; if $S_1+S_2=0$ take $R'=\varnothing$. Put
$\sigma(v)\coloneqq v_1S_2-v_2S_1$, so $\sigma(S)=0$ and
$|\sigma(v_l)|\le\eta(S_1+S_2)$. Order the pieces greedily: if the current
partial sum $p$ has $\sigma(p)\ge0$, append a remaining piece with
$\sigma\le0$, which exists because the remaining sum has $\sigma=-\sigma(p)$;
otherwise append one with $\sigma\ge0$. Then every partial sum satisfies
$|\sigma(p)|\le\eta(S_1+S_2)$. The quantity $p_1+p_2$ increases from $0$ to
$S_1+S_2$ in steps $\le2\eta$, so some partial sum has
$|p_1+p_2-\theta(S_1+S_2)|\le\eta$. Put $w\coloneqq p-\theta S$. Then
\[
 w_1=\frac{\sigma(p)+S_1(w_1+w_2)}{S_1+S_2},\qquad
 w_2=\frac{S_2(w_1+w_2)-\sigma(p)}{S_1+S_2},
\]
so $|w_1|,|w_2|\le2\eta$. Take $\eta=\varepsilon/2$ and let $R'$ be the
union of the pieces in that partial sum.
\end{proof}

\section{One exchange rate}\label{sec:rate}

For disjoint measurable $A,T\subseteq H'$, the \emph{pair} $(A,T)$ is the
change from silence to debt on $A$ and to credit on $T$. Call it \emph{SI}
(a strict improvement) if $V+\debt(A)+\cred(T)\prec_1V$, and \emph{AC}
otherwise. By completeness, AC means $V\preceq_1V+\debt(A)+\cred(T)$. By
(IB) the classification does not change if $V$ is replaced by any regular
model silent on $A\cup T$. Pairwise disjoint pairs that are all SI compose to
an SI pair, and pairs that are all AC compose to an AC pair: apply them one
after another. An SI pair has $T$ non-null, by (G). Put
\[
 r(A,T)\coloneqq P_d(A)/P_c(T)\in[0,\infty],\qquad
 \sigma_-\coloneqq\sup_{\text{SI}}r,\qquad\sigma_+\coloneqq\inf_{\text{AC}}r,
\]
with the convention $r(A,T)\coloneqq\infty$ whenever $P_c(T)=0$, including
the case $P_d(A)=0$.

\begin{lemma}\label{lem:SIAC}
Under the hypotheses of Theorem~\ref{thm:fine}, $0<\sigma_-$ and
$\sigma_+<\infty$.
\end{lemma}
\begin{proof}
Use any edge $[e_1,e_j]$ and Lemma~\ref{lem:approach}. \emph{An SI pair.} Let
$Z\subseteq\Band_j[1,\infty)$ be non-null and $Y\coloneqq E_j(1)+\cred(Z)$,
so $Y\prec_1E_j(1)$. As in Lemma~\ref{lem:debtbands}(a), (T1) gives
$Y\prec_1E_j(s)$ for some $s<1$. The two models differ by the pair
$(\Band_j[s,1),Z)$, which is therefore SI. Its ratio is positive and finite,
by Lemma~\ref{lem:Pd}(ii).

\emph{An AC pair.} Let $A\subseteq\Band_j[0,1)$ be non-null and
$F\coloneqq E^j_1(1)+\debt(A)$, so $F\prec_jE^j_1(1)$. As in
Lemma~\ref{lem:creditbands}(b), (T2) gives $E^j_1(s)\preceq_1F$ for some
$s>1$. So $(A,\Band_j(1,s])$ is AC, with finite ratio.
\end{proof}

\begin{lemma}\label{lem:rates}
Under the hypotheses of Theorem~\ref{thm:fine}, $\sigma_-\le\sigma_+$.
\end{lemma}
\begin{proof}
Suppose $(A,T)$ is AC and $(A_0,T_0)$ is SI with $r(A,T)<r(A_0,T_0)$. Write
$a=P_d(A)$, $t=P_c(T)$, $a_0=P_d(A_0)$, $t_0=P_c(T_0)$. Here $t_0>0$, and
$t>0$ since the ratio of $(A,T)$ is finite. Also $a>0$, because adding credit
alone is a strict improvement, and $a_0>0$, because $a_0/t_0>a/t\ge0$.
Choose integers $K,K_0\ge4$ with
\[
 t_0/t<K_0/K<a_0/a,
\]
which is possible because $t_0/t<a_0/a$. Let $\gamma>0$ be small. By repeated
use of Lemma~\ref{lem:Pd}(iv), split $A$ and $T$ into $K$ pieces each, and
$A_0$ and $T_0$ into $K_0$ pieces each. Each piece carries, for both $P_d$
and $P_c$, within $\gamma$ of $1/K$ (resp.\ $1/K_0$) of the mass of the set
it comes from. Pair the pieces to get $K$ pairs $(A_l,T_l)$ and $K_0$ pairs
$(A_{0l},T_{0l})$.

If all $(A_l,T_l)$ were SI, their composition $(A,T)$ would be SI. So one of
them, say $(A',T')$, is AC. Likewise some $(A_0',T_0')$ is SI. For $\gamma$
small,
\[
 P_d(A')<P_d(A_0'),\qquad P_c(T')>P_c(T_0'),\qquad
 P_c\bigl(H'\setminus(A'\cup A_0')\bigr)\ge1-\tfrac1K-\tfrac1{K_0}-2\gamma
 >P_c(T_0').
\]
The last inequality holds because the left side is at least $\tfrac12-2\gamma$
and $P_c(T_0')\le\tfrac1{K_0}+\gamma\le\tfrac14+\gamma$.
By Lemma~\ref{lem:Pd}(iii) choose $T^*\subseteq H'\setminus(A'\cup A_0')$
with $P_c(T_0')<P_c(T^*)<P_c(T')$. By Lemma~\ref{lem:Pd}(i) and (IB),
\begin{align*}
 V+\debt(A_0')+\cred(T^*)&\prec_1V+\debt(A_0')+\cred(T_0'),\\
 V+\debt(A')+\cred(T^*)&\prec_1V+\debt(A_0')+\cred(T^*),\\
 V+\debt(A')+\cred(T')&\prec_1V+\debt(A')+\cred(T^*).
\end{align*}
The first replaces credit by heavier credit in the context $\debt(A_0')$,
which $T^*$ and $T_0'$ avoid. The second replaces debt by lighter debt in the
context $\cred(T^*)$, which $A'$ and $A_0'$ avoid. The third replaces credit
by heavier credit in the context $\debt(A')$, which $T^*$ and $T'$ avoid.
Overlaps of the two compared sets are allowed, since $\preceq_d,\preceq_c$
compare differences. All models are regular. Hence
$V+\debt(A')+\cred(T')\prec_1V+\debt(A_0')+\cred(T_0')\prec_1V$, so
$(A',T')$ is SI. This is a contradiction.
\end{proof}

Fix $\kappa\in[\sigma_-,\sigma_+]$. Then $0<\kappa<\infty$.

\begin{lemma}\label{lem:kappa}
Let $A,T\subseteq H'$ be disjoint.
\begin{enumerate}[label=(\alph*)]
\item If $P_d(A)<\kappa P_c(T)$, then $(A,T)$ is SI.
\item If $P_d(A)>\kappa P_c(T)$, then $V\prec_1V+\debt(A)+\cred(T)$.
\end{enumerate}
\end{lemma}
\begin{proof}
(a) $r(A,T)<\kappa\le\sigma_+$, so $(A,T)$ is not AC. (b) By
Lemma~\ref{lem:Pd}(iii) choose $A_1\subseteq A$ with
$\kappa P_c(T)<P_d(A_1)<P_d(A)$. Then $A_2\coloneqq A\setminus A_1$ has
positive $P_d$-mass, so it is non-null. Since $r(A_1,T)>\kappa\ge\sigma_-$,
$(A_1,T)$ is not SI, hence AC. Adding debt on $A_2$ is a strict worsening by
strict (G).
\end{proof}

\begin{theorem}\label{thm:extension}
Under the hypotheses of Theorem~\ref{thm:fine}, put
$\Lambda(M)\coloneqq P_d(\debt(M))-\kappa P_c(\cred(M))$ for regular $M$.
Then $M\preceq_1M'$ implies $\Lambda(M)\le\Lambda(M')$ for all regular
$M,M'$.
\end{theorem}
\begin{proof}
Let $\mathcal D\coloneqq\{g\in H':M(g)\ne M'(g)\}$ and let $N$ be $M$
silenced on $\mathcal D$; it is also $M'$ silenced on $\mathcal D$. Let $A,T$
be the debt and credit sets of $M$ inside $\mathcal D$, and $A',T'$ those of
$M'$. Then $M=N+\debt(A)+\cred(T)$ and $M'=N+\debt(A')+\cred(T')$, with
$A\cap T=A'\cap T'=\varnothing$. Put
\[
 \lambda\coloneqq P_d(A)-\kappa P_c(T),\qquad
 \lambda'\coloneqq P_d(A')-\kappa P_c(T'),
\]
so $\Lambda(M)-\Lambda(M')=\lambda-\lambda'$. Suppose $\lambda'<\lambda$; we
show $M'\prec_1M$. Fix $\varepsilon>0$ with $3\varepsilon<\lambda-\lambda'$.
All choices of subsets below use Lemma~\ref{lem:Pd}(iii), and all comparisons
use Lemma~\ref{lem:kappa} in a context silent on the sets changed, via (IB).

\emph{Step 1: from $M$ to a one-signed model $Z\preceq_1M$.}
\begin{itemize}[leftmargin=*,itemsep=1pt]
\item If $\lambda>0$: choose $A^1\subseteq A$ with
$\kappa P_c(T)<P_d(A^1)<\min(\kappa P_c(T)+\varepsilon,P_d(A))$, and put
$Z\coloneqq N+\debt(A\setminus A^1)$. Then $M=Z+\debt(A^1)+\cred(T)$ and
$Z\prec_1M$ by Lemma~\ref{lem:kappa}(b). Here $Z$ has debt only, on a
non-null set, and $l\coloneqq\Lambda(Z)-\Lambda(N)\in(\lambda-\varepsilon,
\lambda)$, $l>0$.
\item If $\lambda\le0$ and $P_d(A)>0$: choose $T^1\subseteq T$ with
$P_d(A)-\varepsilon<\kappa P_c(T^1)<P_d(A)$ and put
$Z\coloneqq N+\cred(T\setminus T^1)$. Then $M=Z+\debt(A)+\cred(T^1)$ and
$Z\prec_1M$ by Lemma~\ref{lem:kappa}(b). Here $Z$ has credit only, and
$l\in(\lambda-\varepsilon,\lambda)$, $l<0$.
\item If $P_d(A)=0$: $A$ is null; put $Z\coloneqq M=N+\cred(T)$, credit
only, with $l=\lambda\le0$.
\end{itemize}

\emph{Step 2: from $M'$ to a one-signed model $Z'$ with $M'\preceq_1Z'$.}
\begin{itemize}[leftmargin=*,itemsep=1pt]
\item If $\lambda'\ge0$ and $P_c(T')>0$: choose $A'^1\subseteq A'$ with
$\kappa P_c(T')-\varepsilon<P_d(A'^1)<\kappa P_c(T')$ and put
$Z'\coloneqq N+\debt(A'\setminus A'^1)$. Then $M'=Z'+\debt(A'^1)+\cred(T')$
and $M'\prec_1Z'$ by Lemma~\ref{lem:kappa}(a). Here $Z'$ has debt only, and
$l'\coloneqq\Lambda(Z')-\Lambda(N)\in(\lambda',\lambda'+\varepsilon)$,
$l'>0$.
\item If $\lambda'<0$: choose $T'^1\subseteq T'$ with
$P_d(A')<\kappa P_c(T'^1)<\min(P_d(A')+\varepsilon,\kappa P_c(T'))$ and put
$Z'\coloneqq N+\cred(T'\setminus T'^1)$. Then $M'=Z'+\debt(A')+\cred(T'^1)$
and $M'\prec_1Z'$ by Lemma~\ref{lem:kappa}(a). Here $Z'$ has credit only, on
a non-null set, and $l'\in(\lambda',\lambda'+\varepsilon)$, $l'<0$.
\item If $\lambda'\ge0$ and $P_c(T')=0$: put $Z'\coloneqq M'=N+\debt(A')$,
debt only, with $l'=\lambda'\ge0$.
\end{itemize}

\emph{Step 3: $Z'\prec_1Z$.} We have $l'<\lambda'+\varepsilon<\lambda-
2\varepsilon<l$. The debt or credit sets of $Z$ lie in $A\cup T$, those of
$Z'$ in $A'\cup T'$; $A\cap A'=T\cap T'=\varnothing$.
\begin{itemize}[leftmargin=*,itemsep=1pt]
\item Both debt only, $Z=N+\debt(B)$, $Z'=N+\debt(B')$: $P_d(B')=l'<l=
P_d(B)$, so $B'\prec_dB$ and $Z'\prec_1Z$.
\item Both credit only, $Z=N+\cred(B)$, $Z'=N+\cred(B')$:
$\kappa P_c(B')=-l'>-l=\kappa P_c(B)$, so $B\prec_cB'$ and $Z'\prec_1Z$.
\item $Z$ debt only and $Z'$ credit only: $Z$ arises only in the first case
of Step~1, on a non-null set, so $N\prec_1Z$; and $Z'\preceq_1N$ by (G).
\item $Z$ credit only and $Z'$ debt only: then $l\le0\le l'$, contradicting
$l'<l$.
\end{itemize}
Every model in the chain $M'\preceq_1Z'\prec_1Z\preceq_1M$ is regular. So
$M'\prec_1M$, and $M\preceq_1M'$ fails.
\end{proof}

\section{Main results}\label{sec:main}

\begin{theorem}\label{thm:A}
Let $n\ge3$, assume (G) at every world and (S), and let world~$i$ be complete
with (IB). If (EP$_{ij}$) holds for every $j\ne i$, then world $i$ has a loss
$L_i$ of the form of Section~\ref{sec:setting}, with \eqref{eq:positive},
such that $D\preceq_iE$ implies $L_i(D)\le L_i(E)$ for all regular $D,E$.
\end{theorem}
\begin{proof}
Relabel so that $i=1$; for $i\ne1$, the losses obtained on the half-space
$\{g_i>0\}$ are transported to $H$ as in the proof of
Theorem~\ref{thm:reduction}. Theorem~\ref{thm:extension} gives $\Lambda$. With $U$
the unanimous problems, put
\[
 \nu_{11}(B)\coloneqq\kappa P_c(B\cap H')+\mu(B\cap U),\qquad
 \nu_{12}(B)\coloneqq P_d(B\cap H')+\mu(B\cap U),\qquad
 \nu_{13}\coloneqq\mu.
\]
These satisfy \eqref{eq:positive} by Lemma~\ref{lem:Pd}(ii). For regular $D$,
the rank of $D$ at world~$1$ is $0$ on $U$, and on $H'$ it is at least $1$ off
$\cred(D)$ and at least $2$ on $\debt(D)$. So
$L_1(D)=\kappa P_c(H'\setminus\cred(D))+P_d(\debt(D))=\kappa+\Lambda(D)$.
\end{proof}

\begin{theorem}\label{thm:reduction}
Let $n\ge3$. Suppose every world's ranking is complete and satisfies (G) and
(IB), and that (S) holds. If (EP$_{ij}$) holds for every edge, then every
coherent model whose credal set contains three vertices is dominated. In
particular $V$ is dominated.
\end{theorem}
\begin{proof}
Theorem~\ref{thm:A} applies at every world. For world $i$ the losses are
built on the representative half-space $\{g_i>0\}$. The map
$g\mapsto\operatorname{sgn}(g_1)\,g$ carries it to $H$ off a null set and
preserves $\mu$, so the losses transport to the form of
Section~\ref{sec:setting}. Lemma~\ref{lem:cboi} finishes.
\end{proof}

For $n=3$ the endpoint property follows from extensions at the other two
worlds. The next proposition is IC's Proposition~3 with finitely additive
losses. Write $\{u,v,w\}=\{1,2,3\}$.

\begin{proposition}\label{prop:FA}
Let $n=3$, assume (G) and (S), let world $u$ satisfy (IB), and suppose $V$ is
undominated. Suppose worlds $v,w$ have losses $L_v,L_w$ of the form of
Section~\ref{sec:setting} with \eqref{eq:positive} and \eqref{eq:W}. Then
every coherent undominated model competitive on $\{u,v\}$ has $e_u$ or $e_v$
in its credal set, and every one competitive on $\{u,w\}$ has $e_u$ or
$e_w$. No completeness is assumed at any world.
\end{proposition}
\begin{proof}
Relabel so that $u=1$, $v=2$, $w=3$, and use IC's notation $P_1,P_2,P_3$ and
$o_K$. Let $C$ be coherent, undominated and competitive on $\{1,2\}$, with
credal interval $[a,b]$ on $[e_1,e_2]$ and $0<a\le b<1$. Then
\[
 T_0\coloneqq P_2\cap\{t_2>(1-a)/a\}\ (o_{13}),\qquad
 B_0\coloneqq P_1\cap\{t_2<(1-b)/b\}\ (o_{23})
\]
are non-null.

By the proof of CBOI's Lemma~3 (Yosida--Hewitt), there are integrable densities
$f_{kl}>0$ a.e.\ and measurable sets $Z_m$ with $\mu(Z_m)\to0$ such that
$\nu_{kl}(Q)=\int_Qf_{kl}\,d\mu$ for every measurable $Q\subseteq H\setminus
Z_m$, for $k=2,3$ and all $l$. Fix $Z\coloneqq Z_m$ with $\mu(Z)$ smaller
than the measures of $T_0$, $B_0$ and $P_3$. If two regular models agree off
a measurable $Q\subseteq H\setminus Z$, finite additivity shows that the
difference of their $L_k$-values is the integral over $Q$ of the difference
of local penalties, with silence penalty $c_k\coloneqq f_{k1}$ and bad
penalty $b_k\coloneqq f_{k1}+f_{k2}$. Thus IC's cost table holds for changes
inside $H\setminus Z$:
\[
 P_1:\ 0\to o_{23}\colon(-c_2,-c_3);\qquad
 P_2:\ 0\to o_{13}\colon(b_2-c_2,-c_3);\qquad
 P_3:\ 0\to o_{12}\colon(-c_2,b_3-c_3).
\]
Now repeat the proof of IC's Proposition~3 inside $T_0\setminus Z$,
$B_0\setminus Z$ and $P_3\setminus Z$. Choose a non-null
$T\subseteq T_0\setminus Z$ with $(b_2-c_2)(T)<c_2(B_0\setminus Z)$, a
non-null $S_*\subseteq P_3\setminus Z$ with $(b_3-c_3)(S_*)<c_3(T)$, and
$\eta>0$ with $2\eta\le c_2(S_*)$ and $\eta<(b_2-c_2)(T)$. The measure
$c_2\,d\mu$ is countably additive and atomless, so there are
$B\subseteq B_0\setminus Z$ with $c_2(B)=(b_2-c_2)(T)-\eta$ and
$S\subseteq S_*$ with $c_2(S)=2\eta$. With $A=T\cup B$ and $X,Y$ as in IC's
Lemma~2,
\[
 L_2(X)-L_2(V)=-\eta,\qquad L_3(X)-L_3(V)<0,\qquad L_2(Y)-L_2(C)=-\eta.
\]
Since $V$ and $C$ are coherent and undominated, \eqref{eq:W} in contrapositive
form gives $V\not\preceq_2X$, $V\not\preceq_3X$ and $C\not\preceq_2Y$. IC's
Lemma~2 contradicts strict (G). The statement for $\{1,3\}$ follows by
exchanging worlds $2$ and $3$.
\end{proof}

\begin{theorem}\label{thm:two}
Let $n=3$ and assume (G) at every world and (S). Suppose two worlds $v,w$
have losses of the form of Section~\ref{sec:setting} with
\eqref{eq:positive} and \eqref{eq:W}, and the third world $u$ has a complete
ranking satisfying (IB). Then $V$ is dominated.
\end{theorem}
\begin{proof}
Suppose $V$ is undominated. The selections at the endpoints of $[e_u,e_v]$
are $D_u$ and $D_v$. The others are coherent, undominated and competitive on
$\{u,v\}$. Likewise on $[e_u,e_w]$. By Proposition~\ref{prop:FA}, (EP$_{uv}$)
and (EP$_{uw}$) hold.
Theorem~\ref{thm:A} gives world $u$ a loss satisfying \eqref{eq:positive} and
the weak-extension property for all regular pairs, hence \eqref{eq:W}.
Lemma~\ref{lem:cboi} makes $V$ dominated, a contradiction.
\end{proof}

\begin{corollary}\label{cor:SH}
Let $n=3$, let worlds $2,3$ be ranked by ordinary additive losses with
strictly ordered local penalties and an integrable envelope, and let world
$1$ have a complete ranking with (G) and (IB). If (S) holds, $V$ is
dominated. In particular no configuration of DT's Section~6 (dense jumps of
$F$ and $G$, separable $\rho$ with unbounded jump ratios) can occur with a
complete world~$1$.
\end{corollary}

\begin{corollary}\label{cor:shape}
Let $n=3$, let all rankings be complete with (G) and (IB), assume (S), and
suppose $V$ is undominated. Then:
\begin{enumerate}[label=(\alph*)]
\item at most one world has a loss of the form of
Section~\ref{sec:setting} with \eqref{eq:positive} and \eqref{eq:W}; by EDL's
Theorem~2, at least two worlds carry leverage sequences;
\item some edge has a selection whose credal set contains neither endpoint.
\end{enumerate}
\end{corollary}
\begin{proof}
(a) is Theorem~\ref{thm:two}, applied to any two worlds with such losses.
(b) is Theorem~\ref{thm:reduction}.
\end{proof}

\section{Remarks and what remains}\label{sec:remarks}

\begin{remark}\label{rem:complete}
\emph{Completeness.} Completeness of the free world is used in (T1), in the
comparative-probability step (Lemma~\ref{lem:cp} and
Theorem~\ref{thm:SW}), and in the SI/AC dichotomy. (T2),
Lemma~\ref{lem:edge}, Proposition~\ref{prop:FA} and Lemma~\ref{lem:cboi} do
not use it. For an incomplete free world, (T1) gives only non-comparisons
$C_q\not\preceq_1Y$. A variant of (S) that also requires the strict sections
$\{p:X\prec_iC_p\}$ to be open restores (T1). For complete rankings this
variant is equivalent to (S). We have not checked whether the rest of the
argument then goes through for incomplete rankings; Wakker's theorem needs a
complete order. The sub-case of IC and DT with an \emph{incomplete} world~$1$
therefore remains open. For it, the chain results of IC and DT, which never
use completeness, are the strongest known.
\end{remark}

\begin{remark}\label{rem:DT}
\emph{Status of DT.} A referee checked every result of DT.
\begin{itemize}[leftmargin=*,itemsep=2pt]
\item Lemma~1 (Bayes models $B_\theta$), Theorem~2 (dense thresholds), and
Corollaries~3 and~4 are correct. Theorem~2 needs three small additions: cite
the face lemma for the full-support step; use IC's weights lemma, part (c),
for full-support indices as well as for indices on $(e_2,e_3)$; and state
that the limit $C_*$ minimizes $L_2+\kappa L_3$.
\item Lemma~5 (debt contraction) has a gap. The silence of a good base on
$P_1$ lies in $\{t_{13}>\sigma\}$, where
$\sigma\coloneqq\inf\{m_3/m_1:m\in K,\ m_1>0\}$. Nothing bounds $\sigma$, so
the reservoir $M_0\subseteq\{t_{13}<\theta_0\}$ with $2\theta_0<h$ need not
exist. Theorem~6 (one-strip bound) and Corollary~7 inherit the gap.
\item \emph{Repair.} Replace step~2 of Lemma~5 by the change
$[Y\colon o_2\to0,\ Y'\colon0\to o_{12},\ M\colon0\to o_{23}]$ at the good
base, with $u(Y')$ chosen freely. Let $a\coloneqq\operatorname{ess\,sup}_YF$
and $b\coloneqq\operatorname{ess\,inf}_{Y'_0}F$. The change is accepted
whenever $c_3(M)+c_2(M)/b>(1/b-1/a)\,c_2(Y)$, for any $M$ in the silence;
no bound on $\rho$ over $M$ is needed. This gives: if $\rho\ge A>0$ on
$P_1\cap\{t_{12}<h\}$ and some good base has
$1/F(\tau-)-1/F(\tau+)<A$, then $V$ is dominated. With Theorem~2 at a
continuity point of $F$, and its mirror, this yields a stronger form of
Theorem~6: $V$ is dominated if $\operatorname{ess\,inf}\rho>0$ on some
$P_1\cap\{t_{12}<h\}$, or $\operatorname{ess\,sup}\rho<\infty$ on some
$P_1\cap\{t_{13}<h\}$. Corollary~7 and the strip claims of Section~6 then
hold.
\item Lemma~8 and Theorem~9 (separable ratios) are correct; Lemma~8 even
holds with $\lambda=F(\tau-)/\sup F$. Section~6 should add the mirror
statement for $h$.
\end{itemize}
A further refereed fact for the sub-case of IC: if $\rho$ is essentially
unbounded above on $P_1\cap\{t_{12}<\ell\}$ for \emph{some} $\ell$, then $V$
is dominated, directly by IC's ratio criterion and one $E_2(\ell)$ move; no
good base is needed. This strengthens IC's Theorem~17(a). These results are
superseded by Corollary~\ref{cor:SH} when world $1$ is complete. They remain relevant for an incomplete world~$1$. So does
the separability theorem for $\rho=c_3/c_2$ and
$\sigma=(b_3-c_3)/(b_2-c_2)$ on $P_1$, which is written up in full
separately.
\end{remark}

\begin{remark}\label{rem:open}
\emph{What remains.} For complete rankings and any $n\ge3$, the conjecture
now reduces to the following statement: if $V$ is undominated, then (EP)
holds on every edge. Indeed, if (EP) holds everywhere,
Theorem~\ref{thm:reduction} makes $V$ dominated. The endpoint property on the
two edges at a world is known when the other two worlds have extensions
(Proposition~\ref{prop:FA}), and, for $n=3$, when they satisfy the
path-closedness at $V$ and at the models $V+o_{12}(S)$ used in IC's
Theorem~5. We could not prove it in general. By
Corollary~\ref{cor:shape}, a counterexample with complete rankings needs an
edge with a two-sided selection and leverage at two worlds. Neither is
excluded by the arguments above.
\end{remark}

\begin{thebibliography}{9}
\bibitem{Savage} L.~J.~Savage, \emph{The Foundations of Statistics}, Wiley,
New York, 1954. Chapter~3, \S3.
\bibitem{Wakker} P.~Wakker, Agreeing probability measures for comparative
probability structures, \emph{Annals of Statistics} 9 (1981), 658--662.
doi:10.1214/aos/1176345469.
\end{thebibliography}
\end{document}
